\documentclass[conference]{IEEEtran}

\usepackage{amsmath, amssymb, amsfonts}
\usepackage{graphicx}
\usepackage{cite}
\usepackage{booktabs}
\usepackage{hyperref}
\usepackage{url}
\usepackage{algorithm}
\usepackage{algpseudocode}
\usepackage{tikz}
\usetikzlibrary{positioning, shapes, arrows.meta, calc, fit, backgrounds}
\usepackage{fontawesome5} % For the icons

\title{Sponge Attacks Against AI-Based Models in the Smart Grid: Analysis and Experimentation}

\author{\IEEEauthorblockN{Bora Pilav}
\IEEEauthorblockA{Institute for Automation and Applied Informatics (IAI)\\
Karlsruhe Institute of Technology\\
Karlsruhe, Germany\\
Email: ugglf@student.kit.edu}
}

\markboth{Seminar on Adversarial Machine Learning}{Sponge Attacks in Smart Grids}

\begin{document}
\maketitle

\begin{abstract}
The integration of Artificial Intelligence (AI) into the Smart Grid is critical for optimizing energy distribution, forecasting demand, and detecting anomalies in real-time. However, the deployment of these complex models introduces significant security vulnerabilities. While traditional adversarial attacks target model integrity to induce misclassifications, this report investigates ``Sponge Attacks''—a class of denial-of-service (DoS) attacks targeting model availability. Sponge attacks exploit the computational density of deep learning models to exhaust system resources, specifically maximizing inference latency and energy consumption. In the context of the Smart Grid, such delays can destabilize grid operations or drain battery-powered edge devices like smart meters. This paper conducts a meta-analysis of existing literature on sponge examples and poisoning, identifies vulnerable components within smart grid architectures, and presents a proof-of-concept experiment demonstrating the impact of these attacks on model performance. Finally, we explore mitigation strategies to secure grid AI systems against resource exhaustion.
\end{abstract}

\begin{IEEEkeywords}
Smart Grid, Sponge Attacks, Adversarial Machine Learning, Energy-Latency, Availability Attacks, Deep Learning Security.
\end{IEEEkeywords}

\IEEEpeerreviewmaketitle

\section{Introduction}
The modernization of the electrical power infrastructure into the ``Smart Grid'' represents a paradigm shift from centralized, one-way transmission to a decentralized, bi-directional network driven by data. To manage the increasing complexity of renewable energy integration and dynamic load balancing, operators are rapidly adopting Artificial Intelligence (AI) and Machine Learning (ML) techniques. These models are deployed for critical tasks such as load forecasting, fault detection, and real-time pricing optimization \cite{sanchez2024attacking}.

Despite their operational benefits, the reliance on deep learning models expands the grid's attack surface. Existing research in adversarial machine learning has predominantly focused on \textit{integrity} attacks, where an adversary perturbs input data (e.g., adversarial examples) to cause incorrect model predictions. However, in safety-critical systems like the Smart Grid, \textit{availability} is equally paramount. Decisions regarding frequency regulation and fault isolation must be made within milliseconds to prevent cascading failures.

This report addresses a specific threat to availability known as the ``Sponge Attack,'' formally introduced by Shumailov et al. \cite{shumailov2021sponge}. Unlike standard adversarial attacks that seek to fool a model, sponge attacks aim to ``soak up'' computational resources. By crafting specific inputs (Sponge Examples) \cite{cina2022sponge} or poisoning model weights during training (Sponge Poisoning) \cite{cina2022sponge}, an attacker can trigger the worst-case computational performance of a neural network. This results in significant spikes in inference latency and energy consumption.

The implications for the Smart Grid are severe. A sponge attack targeting a central load forecasting server could introduce latency that desynchronizes grid control loops. Furthermore, attacks on edge devices—such as Advanced Metering Infrastructure (AMI) or IoT sensors—can rapidly drain batteries, rendering the monitoring infrastructure blind.

To support reproducible research in Smart Grid security, we release our code, pre-trained ACT-LSTM models, and the XAI-guided attack pipeline at a public repository\footnote{\url{https://github.com/student-kit/sponge-grid-attacks}}.

This paper is structured to address three primary objectives:
\begin{itemize}
    \item \textbf{Literature Review:} A meta-analysis of AI-based smart grid systems and the mechanics of sponge attacks (Section \ref{sec:related}).
    \item \textbf{System Model:} A formal definition of the threat landscape and target architectures (Section \ref{sec:threat_model}).
    \item \textbf{Experimentation:} A rigorous comparison of Black-Box vs. White-Box attacks and an analysis of defense mechanisms (Section \ref{sec:experiments} and \ref{sec:results}).
\end{itemize}

\section{Related Work}
\label{sec:related}

This section reviews the current state of Artificial Intelligence (AI) integration within the Smart Grid, follows with an analysis of the fundamental mechanics of Sponge Attacks, and concludes with recent advancements in Sponge Poisoning and Moving Target Defense (MTD).

\subsection{AI in Smart Grid Systems}
The transition to the Smart Grid has necessitated the deployment of learning-based components to handle the stochastic nature of renewable energy generation and variable consumer demand. S{\'a}nchez et al. \cite{sanchez2024attacking} highlight that these models are ubiquitous, appearing in critical use cases such as energy demand prediction, anomaly detection for cybersecurity, and preventative maintenance. 

For instance, Recurrent Neural Networks (RNNs) and Long Short-Term Memory (LSTM) networks are standard for time-series load forecasting due to their ability to capture temporal dependencies. Simultaneously, Convolutional Neural Networks (CNNs) \cite{sanchez2024attacking, shapira2023phantom} are increasingly used for visual inspection of grid infrastructure (e.g., drone-based power line monitoring). However, the security of these models is often overlooked. While most research focuses on \textit{integrity attacks}—where an adversary manipulates smart meter data to lower bills or hide consumption spikes—\textit{availability attacks} that target the underlying hardware resources remain under-explored in the specific context of grid infrastructure.

\subsection{Sponge Examples: Inference-Time Attacks}
The concept of Sponge Examples was formally introduced by Shumailov et al. \cite{shumailov2021sponge} as a novel Denial-of-Service (DoS) vector against machine learning models. Unlike adversarial examples that aim for misclassification, sponge examples are optimized to maximize the energy consumption and inference latency of a model.

Shumailov et al. identify two primary dimensions of computation that adversaries can exploit:
\begin{itemize}
    \item \textbf{Computation Depth:} For models with dynamic execution paths (e.g., language models or RNNs used in sequencing), an attacker can craft inputs that force the model to process the maximum number of tokens or steps, preventing early-exit optimizations.
    \item \textbf{Data Sparsity:} Many hardware accelerators (GPUs/TPUs) utilize optimizations like zero-skipping to avoid performing multiplications with zero-valued inputs. Sponge examples are crafted to be dense (few zeros), thereby disabling these hardware optimizations and increasing the arithmetic load.
\end{itemize}

In the domain of computer vision, which is relevant for grid surveillance, Shapira et al. \cite{shapira2023phantom} demonstrated attacks against object detectors. They exploited the Non-Maximum Suppression (NMS) algorithm—a post-processing step used to filter overlapping bounding boxes. By generating "Phantom Sponges" that trigger a worst-case number of candidate boxes, they successfully increased inference latency without altering the visual content significantly.

\subsection{Explainable AI (XAI) as an Attack Vector}
Explainable AI (XAI) methods, such as SHAP and Integrated Gradients (IG), were originally designed to build trust by highlighting which input features contribute to a model's prediction. However, recent research has demonstrated that these same tools can be weaponized by adversaries to identify model vulnerabilities.

\subsubsection{XAI for Integrity Attacks}
S{\'a}nchez et al. \cite{sanchez2024attacking} utilized feature attribution maps to craft more potent adversarial examples against smart grid intrusion detection systems. By perturbing only the features with the highest Gini importance, they achieved high evasion rates with minimal input distortion.

\subsubsection{XAI for Model Stealing}
In a parallel domain, \cite{steganographic2026} demonstrated that XAI can facilitate model extraction. By analyzing the gradient maps of a target network, attackers can infer the architecture's depth and kernel size. Our work extends this "weaponized XAI" paradigm to availability attacks: we use IG not to interpret the model's \textit{output}, but to interpret its \textit{execution time}, effectively reverse-engineering the logical trigger conditions of the halting unit.

% --- TABLE I: TAXONOMY ---
\begin{table*}[t]
    \centering
    \caption{Taxonomy of Sponge Attacks and their Potential Impact on Smart Grid Infrastructure.}
    \label{tab:literature_review}
    \renewcommand{\arraystretch}{1.3}
    \begin{tabular}{l p{3cm} p{3.5cm} p{4cm} p{3cm}}
        \toprule
        \textbf{Reference} & \textbf{Attack Technique} & \textbf{Target Mechanism} & \textbf{Smart Grid Vulnerability} & \textbf{Grid Impact} \\
        \midrule
        Shumailov et al. \cite{shumailov2021sponge} & Sponge Examples (Inference) & Computation Depth (NLP) \& Sparsity (Vision) & Centralized Forecasting Servers (LSTMs/Transformers) & Delayed dispatch decisions; Real-time pricing lag. \\
        \hline
        Shapira et al. \cite{shapira2023phantom} & Phantom Sponges & Non-Maximum Suppression (Object Detection) & Autonomous Inspection Drones (YOLO/R-CNN) & Battery exhaustion; Unfinished inspection routes. \\
        \hline
        Cinà et al. \cite{cina2022sponge} & Sponge Poisoning & Training-time Loss Manipulation & Edge Devices (Smart Meters) & Hardware thermal throttling; Denial of Service (DoS). \\
        \hline
        Te Lintelo et al. \cite{telintelo2024skip} & SkipSponge & Weight Poisoning (Reducing Sparsity) & Resource-Constrained IoT Sensors & Increased base-load energy consumption. \\
        \hline
        \textbf{This Work} & \textbf{Black-Box Genetic Algo} & \textbf{Dynamic Control Flow (Halting Unit)} & \textbf{Adaptive Forecasting (ACT-LSTM)} & \textbf{1.6x Latency Spike; Missed aggregation windows.} \\
        \bottomrule
    \end{tabular}
\end{table*}

\subsection{Sponge Poisoning: Training-Time Attacks}
While standard sponge examples require altering the input at inference time, recent research has explored Sponge Poisoning, where the model itself is compromised during the training or update phase.

Cinà et al. \cite{cina2022sponge} explored this in the context of edge devices, which is highly relevant for smart meters and IoT sensors. They demonstrated that by injecting a small percentage of poisoned data into the training set, an attacker can alter the model's objective function to prefer computationally expensive paths. 

More recently, Te Lintelo et al. \cite{telintelo2024skip} introduced \textit{SkipSponge}, a technique that targets the weights of pre-trained models directly. This method is particularly stealthy as it requires fewer data samples than traditional poisoning. By manipulating the weights to reduce the sparsity of intermediate activations, SkipSponge forces the hardware to perform more calculations, effectively turning a standard efficient model into an energy drain.

\subsection{Sensing AI and IoT Vulnerabilities}
The impact of sponge attacks is most severe on battery-constrained devices. Recent work in 2025 has shifted focus from cloud servers to "Sensing AI" (wearable and IoT sensors). 

Hasan et al. \cite{hasan2025sensing} conducted a systematic study on wearable AI. They found that sponge attacks on these devices lead to rapid battery depletion and thermal throttling, effectively taking the sensor offline. This is directly analogous to the risk facing Smart Meters (AMI) in the grid. Their work suggests that standard model compression techniques, such as \textbf{Model Pruning}, can inadvertently serve as a defense by reducing the overall capacity available to be "sponged," though this comes at an accuracy trade-off.

\subsection{Mitigation Strategies}
Defending against availability attacks requires different paradigms than integrity defenses.

\subsubsection{Static Defenses}
Proposed mitigations include \textbf{Fine-Pruning} \cite{telintelo2024skip}, which attempts to "heal" a poisoned model by pruning and re-training affected weights, and \textbf{Input Filtering}, which rejects samples with abnormal energy profiles. However, filters can often be bypassed by adaptive adversaries who optimize for both latency and stealth (as we demonstrate in Section \ref{sec:experiments}).

\subsubsection{Moving Target Defense (MTD)}
A promising direction in Smart Grid security is Moving Target Defense (MTD), which dynamically randomizes system configurations to invalidate attacker knowledge. S{\'a}nchez et al. \cite{sanchez2025lightweight} proposed a lightweight MTD that randomizes Modbus TCP slave addresses to prevent network intrusion, maintaining 95\% detection accuracy.

However, we argue that \textbf{Network-level MTD is insufficient for Model-level availability}. While randomizing IP addresses protects the transport layer, it does not protect the compute layer once a request is authenticated. If the underlying forecasting model remains static, an authenticated adversary can still inject sponge inputs. A true "AI-MTD" would require randomizing the neural architecture itself (e.g., switching between ensemble members with different halting policies).

\section{System Overview and Threat Model}
\label{sec:threat_model}

We formally define the Smart Grid forecasting environment, the specific dynamic architectures targeted, and the adversarial capabilities assumed in our worst-case analysis.

\subsection{Smart Grid Control Loop}
Consider a grid control node (e.g., a frequency regulator) that operates in discrete time steps $T$. At each step, the system receives a multivariate time-series input $X_t \in \mathbb{R}^{F \times W}$ (where $F$ is the number of features and $W$ is the history window) and relies on a deep learning model $f_\theta$ to predict the power generation $Y_{t+1}$.
\begin{equation}
    Y_{t+1} = f_\theta(X_t)
\end{equation}
Crucially, the control action $u_t$ must be computed and applied within a strict dispatch window $\tau_{max}$ (e.g., 300ms for primary frequency response). If the inference latency $L(f_\theta(X_t)) > \tau_{max}$, the control loop opens, leading to uncompensated frequency deviation.

\subsection{Target Dynamics: ACT-LSTM}
We focus on the Adaptive Computation Time (ACT) LSTM, which introduces dynamic depth. For an input $x_t$, the model executes an internal recurrence loop for $N(x_t)$ steps. The halting condition is governed by a sigmoidal halting unit $h$:
\begin{equation}
    N(x_t) = \min \{ n : \sum_{k=1}^{n} h(s_k) \ge 1 - \epsilon \}
\end{equation}
where $s_k$ is the internal state at ponder step $k$. This creates a variable computational cost $C(x) \propto N(x)$. Our specific target is to find an input $x^*$ that maximizes $N(x^*)$ up to a hard limit $M$.

\subsection{Reference Architectures (Static Baselines)}
To validate the unique vulnerability of dynamic compute, we contrast the ACT-LSTM against two static baselines where the computational cost $C(x)$ is constant:
\begin{itemize}
    \item \textbf{DeepAR (RNN):} A standard stacked LSTM where the unrolling depth is fixed to the input window size $W$. The inference time is $O(W)$ regardless of input complexity.
    \item \textbf{Chronos-v2 (Transformer):} A foundation model that utilizes input quantization and fixed-length attention masks. Despite the theoretical $O(W^2)$ complexity of attention, the implementation enforces a deterministic execution path, rendering $L(f_\theta(x)) \approx \text{const}$.
\end{itemize}

\subsection{Threat Model Formulation}
We adopt a comprehensive threat model considering both black-box and white-box capabilities to establish lower and upper bounds on vulnerability.

\subsubsection{Adversary Goals ($\mathcal{G}$)}
The adversary $\mathcal{A}$ aims to violate the availability property of the grid.
\begin{itemize}
    \item \textbf{Primary Goal:} Maximize Latency. Find $\delta$ such that $L(f_\theta(X + \delta)) > \tau_{max}$.
    \item \textbf{Secondary Goal:} Maximize Energy. Maximize $\int P(t) dt$ to drain battery-powered edge sensors.
    \item \textbf{Constraint:} Stealthiness. The perturbation must remain undetectable: $||\delta||_p < \epsilon_{noise}$ and the input must satisfy physical constraints (e.g., Wind Speed $\ge 0$).
\end{itemize}

\subsubsection{Adversary Knowledge ($\mathcal{K}$)}
We define two distinct threat profiles:
\begin{itemize}
    \item $\mathcal{K}_{BB}$ (Black-Box): The attacker knows the architecture (ACT-LSTM) and input format but has no access to weights $\theta$ or gradients. They rely on query-access to measure latency.
    \item $\mathcal{K}_{WB}$ (White-Box): The attacker has full access to $\theta$ and $\nabla_\theta$. This represents a "Worst-Case" scenario (e.g., an insider threat or leaked model) used to calculate the theoretical maximum latency.
\end{itemize}

\section{Methodology}
\label{sec:experiments}

This section details the experimental setup designed to validate the vulnerability of modern time-series forecasters to both black-box and white-box sponge attacks.

\subsection{Dataset and Preprocessing}
We utilize the \textbf{Wind Power Generation Data} dataset acquired from Kaggle, which provides high-resolution time-series data for renewable energy forecasting.
\begin{itemize}
    \item \textbf{Features:} The dataset includes critical meteorological and operational features such as Wind Power (MW), Wind Speed (m/s), Wind Direction ($^\circ$), and ambient temperature.
    \item \textbf{Normalization:} To ensure stable gradient descent during the attack generation, all continuous features were scaled using Z-score normalization.
    \item \textbf{Input Window:} We structured the data into sliding windows of length $T=64$ hours to predict the subsequent $T=10$ hours, enabling the model to capture diurnal cycles and weather fronts.
\end{itemize}

\subsection{Target Architectures}
We contest three distinct architectures to test our hypothesis regarding Static versus Dynamic computation graphs.
\begin{itemize}
    \item \textbf{DeepAR:} A standard 2-layer LSTM (Hidden Size: 128) representing static graphs.
    \item \textbf{Chronos-v2:} A pretrained T5-based foundation model representing quantized, deterministic graphs.
    \item \textbf{ACT-LSTM:} A custom implementation of Adaptive Computation Time, representing dynamic graphs with vulnerable control flow.
\end{itemize}

\subsection{Attack Implementation}
To evaluate the full spectrum of risk defined in our threat model, we implement two distinct attack vectors.

\subsubsection{Baseline: Genetic Algorithm (Black-Box)}

\begin{figure}[t]
\centering
\begin{tikzpicture}[
    node distance=1.2cm and 1.5cm,
    every node/.style={font=\sffamily\footnotesize, align=center},
    % Styles
    block/.style={rectangle, draw=black!80, thick, rounded corners, minimum width=2cm, minimum height=1cm, fill=white, drop shadow},
    process/.style={rectangle, draw=blue!80, thick, rounded corners, minimum width=2.5cm, minimum height=1.2cm, fill=blue!5},
    cycle/.style={circle, draw=green!60!black, thick, fill=green!5, minimum size=2.5cm},
    arrow/.style={-Latex, thick, draw=black!70}
]

    % --- 1. The Target (Opaque) ---
    \node[block, fill=black!80, text=white] (model) {\textbf{Target Model} \\ (Black Box) \\ \tiny Unknown Weights $\theta$};

    % --- 2. The Interaction ---
    % Input Population
    \node[block, left=of model, label=above:\textbf{Population $P_t$}] (pop) {
        \faFileWaveform \ \faFileWaveform \\[-0.3em] \tiny Candidate Inputs
    };
    
    % Output Feedback
    \node[right=of model, text width=2cm] (feedback) {Latency $L(x)$ \\ \tiny (Fitness Score)};

    % --- 3. The Genetic Engine (Evolutionary Cycle) ---
    \node[process, below=1.5cm of model] (ga) {\textbf{Genetic Algorithm} \\ \tiny Selection $\rightarrow$ Crossover $\rightarrow$ Mutation};

    % --- CONNECTIONS ---
    
    % Forward Pass (Query)
    \draw[arrow] (pop) -- node[above, font=\tiny] {Query} (model);
    \draw[arrow] (model) -- (feedback);
    
    % Feedback Loop
    \draw[arrow] (feedback.south) |- (ga.east);
    \draw[arrow] (ga.west) -| (pop.south);
    
    % Annotations
    \node[below=0.1cm of ga, font=\tiny, color=black!60] {Repeat $N$ Generations};

\end{tikzpicture}
\caption{\textbf{Black-Box Threat Model (Genetic Algorithm).} The adversary has no access to model internals. They maintain a population of input waveforms, query the black-box model to measure execution time, and evolve the population to maximize latency.}
\label{fig:blackbox_attack}
\end{figure}

As described in previous work \cite{shumailov2021sponge}, we utilize a Genetic Algorithm (PyGAD) to evolve adversarial inputs. The fitness function is defined solely by the wall-clock inference time: $F(x) = -\text{Latency}(x)$. This assumes $\mathcal{K}_{BB}$.

\subsubsection{Proposed: XAI-Guided Saliency Attack (White-Box)}
\begin{figure}[t]
\centering
\begin{tikzpicture}[
    node distance=1.2cm and 1.5cm,
    every node/.style={font=\sffamily\footnotesize, align=center},
    % Styles
    block/.style={rectangle, draw=black!80, thick, rounded corners, minimum width=2cm, minimum height=1cm, fill=white, drop shadow},
    whitebox/.style={rectangle, draw=orange!80, thick, rounded corners, minimum width=2.5cm, minimum height=1.5cm, fill=white},
    arrow/.style={-Latex, thick, draw=black!70},
    gradient/.style={draw=red!80, thick, dashed, color=red}
]

    % --- 1. The Target (Transparent) ---
    \node[whitebox] (model) {};
    % Internals visible
    \node[font=\tiny, anchor=north] at (model.north) {\textbf{Target Model} (White Box)};
    \node[draw, circle, inner sep=1pt, minimum size=0.5cm, fill=orange!10] at (model.center) (loop) {$\circlearrowright$};
    \node[below=0.1cm of loop, font=\tiny] {Ponder Loop $h(s_t)$};

    % --- 2. The Flow ---
    \node[block, left=of model] (input) {$X + \delta$};
    \node[right=of model, text width=2.5cm] (loss) {\textbf{Ponder Loss} \\ $\mathcal{L} = \sum p_t$};

    % --- 3. The Adversary (Gradient Step) ---
    \node[block, below=1.5cm of model, fill=red!5, draw=red!50] (attacker) {
        \textbf{PGD Update} \\ 
        \tiny $x \leftarrow x + \alpha \cdot \text{sign}(\nabla_x \mathcal{L})$
    };

    % --- CONNECTIONS ---
    
    % Forward Pass
    \draw[arrow] (input) -- (model);
    \draw[arrow] (model) -- (loss);
    
    % Gradient Backprop (The Attack)
    \draw[arrow, gradient] (loss.south) |- node[right, font=\tiny, pos=0.2] {Backprop $\nabla_\theta$} (attacker.east);
    \draw[arrow, gradient] (attacker.west) -| node[left, font=\tiny, pos=0.2] {Perturb $\delta$} (input.south);

    % Visuals
    \node[above=0.1cm of input] {\faSyringe}; % Injection icon

\end{tikzpicture}
\caption{\textbf{White-Box Threat Model (Saliency-Guided PGD).} The adversary inspects the internal halting probabilities (Ponder Loss). By computing the gradient of computation time with respect to the input ($\nabla_x \mathcal{L}$), they iteratively optimize the perturbation $\delta$ to maximize the loop duration.}
\label{fig:whitebox_attack}
\end{figure}

For the $\mathcal{K}_{WB}$ scenario, we propose a novel \textit{Saliency-Guided Sponge Attack}. Instead of maximizing classification error, we define a "Ponder Loss" $\mathcal{L}_{ponder}$ representing the total accumulated halting probability:
\begin{equation}
    \mathcal{L}_{ponder}(x) = \sum_{t=1}^{T} \sum_{k=1}^{M} h(s_{t,k})
\end{equation}
We compute the gradient $\nabla_x \mathcal{L}_{ponder}$ to generate a Saliency Map. This map identifies the specific time-steps that trigger the halting unit. We then apply Projected Gradient Descent (PGD) masked by this saliency map, effectively "pushing" the sensitive time-steps towards the decision boundary where the model hesitates the most.

\subsection{Formal Algorithm: Saliency-Guided PGD}
To ensure reproducibility, we formalize our White-Box attack procedure in Algorithm \ref{alg:sponge}. Unlike standard PGD which maximizes classification loss $\mathcal{L}_{CE}$, our algorithm maximizes the Ponder Cost $\mathcal{L}_{ponder}$ (Eq. 5) while respecting the physical constraints of the wind sensor data (e.g., wind speed cannot be negative).

\begin{algorithm}[h]
\caption{Saliency-Guided Sponge Attack (PGD)}
\label{alg:sponge}
\begin{algorithmic}[1]
\Require Model $f_\theta$, Input $x$, Step size $\alpha$, Iterations $K$, Max Perturbation $\epsilon$
\Ensure Sponge Example $x_{adv}$

\State $x_{adv} \leftarrow x$
\State $x_{orig} \leftarrow x$

\For{$k = 1$ to $K$}
    \State \textit{// Forward pass to get Halt Probabilities}
    \State $N, \{h_t\} \leftarrow f_\theta(x_{adv})$
    
    \State \textit{// Compute Ponder Loss (Sum of Halting Probs)}
    \State $\mathcal{L}_{ponder} \leftarrow \sum_{t} \sum_{step} h_t(step)$
    
    \State \textit{// Compute Gradient w.r.t Input}
    \State $\nabla_x \mathcal{L} \leftarrow \frac{\partial \mathcal{L}_{ponder}}{\partial x_{adv}}$
    
    \State \textit{// Compute Saliency Mask (Top 10\% features)}
    \State $M \leftarrow \mathbb{I}(|\nabla_x \mathcal{L}| > \text{percentile}(|\nabla_x \mathcal{L}|, 90))$
    
    \State \textit{// Update with Masked Gradient Ascent}
    \State $x_{adv} \leftarrow x_{adv} + \alpha \cdot \text{sign}(\nabla_x \mathcal{L}) \cdot M$
    
    \State \textit{// Project back to $\epsilon$-ball and valid range}
    \State $\delta \leftarrow \text{clip}(x_{adv} - x_{orig}, -\epsilon, \epsilon)$
    \State $x_{adv} \leftarrow \text{clip}(x_{orig} + \delta, 0, \text{max\_val})$
\EndFor

\State \Return $x_{adv}$
\end{algorithmic}
\end{algorithm}

The key innovation in Algorithm \ref{alg:sponge} is the masking step (Line 9). By filtering the gradient update $M$, we ensure that perturbations are applied \textit{only} to the time-steps that strongly influence the halting unit. This prevents the "high-frequency noise" look of standard adversarial examples, making our sponge inputs visually smoother and harder to detect.

\section{Experimental Results}
\label{sec:results}

We evaluated the vulnerability of the target architectures under both Black-Box (Genetic Algorithm) and White-Box (Saliency-Guided PGD) threat models. The results quantify the "Sponge Effect" in terms of computational depth, energy consumption, and transferability.

\subsection{Attack Efficacy: White-Box vs. Black-Box}
A key contribution of this work is the comparative analysis between derivative-free optimization and gradient-based attacks. Figure \ref{fig:ponder_vis} (Top Right) illustrates the superior efficacy of the white-box approach.

\subsubsection{Computational Depth}
The ACT-LSTM baseline processes benign inputs with an average computational depth of $\mu=4.61$ ponder steps per time-step. 
\begin{itemize}
    \item \textbf{Genetic Algorithm (GA):} The black-box attack increased this average to $6.45$ steps ($+40\%$).
    \item \textbf{PGD (Ours):} The gradient-based PGD attack achieved an average of $7.64$ steps ($+66\%$), significantly outperforming the GA.
\end{itemize}
As shown in Figure \ref{fig:ponder_vis} (Top Left), the PGD attack forces the model to hit the hard cap of 20 ponder steps at specific intervals (e.g., $t \approx 55$), a behavior rarely achieved by the GA. The histogram (Bottom Right) confirms this, showing a "heavy tail" distribution for PGD that is absent in the baseline.

% --- TABLE II: COMPARISON ---
\begin{table}[t]
    \centering
    \caption{Comparison of Attack Efficacy on ACT-LSTM. The White-Box (PGD) attack significantly outperforms the Black-Box (GA) baseline in all metrics.}
    \label{tab:attack_comparison}
    \renewcommand{\arraystretch}{1.2}
    \resizebox{\columnwidth}{!}{%
    \begin{tabular}{l c c c c}
        \toprule
        \textbf{Method} & \textbf{Ponder Depth} & \textbf{Latency (ms)} & \textbf{Energy (mJ)}\\
        \midrule
        \textbf{Baseline (Benign)} & 4.61 & 47.00 & 696.16\\
        \midrule
        \textbf{Black-Box (GA)} & 6.45 & $\approx 58.0$ (\textbf{+23\%})  & $\approx 850$ (\textbf{+22\%}) \\
        \textbf{White-Box (PGD)} & \textbf{7.64} & \textbf{70.69} (\textbf{+50\%}) & \textbf{1110.70} (\textbf{+68\%}) \\
        \bottomrule
    \end{tabular}%
    }
\end{table}

\begin{figure*}[t]
    \centering
    \includegraphics[width=\textwidth]{act_ponder_visualization.png}
    \caption{\textbf{Dynamics of the Sponge Attack.} (Top Left) Ponder steps per time-increment. The PGD attack (orange) forces the model to "think" for the maximum 20 steps at specific intervals. (Top Right) The White-Box PGD attack outperforms the Black-Box GA, inducing higher average computation (7.64 vs 6.50). (Bottom Right) Histogram showing the PGD attack successfully pushing the model into the "Hard Limit" zone.}
    \label{fig:ponder_vis}
\end{figure*}

\subsection{Energy and Latency Impact}
The increase in computational depth translates directly to physical resource exhaustion.
\begin{itemize}
    \item \textbf{Latency:} Under the PGD attack, the per-sample latency increased from $47.00$ ms to $70.69$ ms ($+50.4\%$) (see Fig. \ref{fig:latency_res}).
    \item \textbf{Energy:} The total energy consumption per inference spiked from $696.16$ mJ to $1110.70$ mJ ($+59.5\%$) (see Fig. \ref{fig:energy_res}). 
\end{itemize}
Crucially, the Power draw remained constant at $\approx 13.5$ W. This confirms that the "Sponge" effect is purely temporal; the attack does not increase the instantaneous intensity of the computation but rather prolongs its duration, draining battery-constrained devices over time.

\begin{figure}[h]
    \centering
    \includegraphics[width=\columnwidth]{act_pgd_latency_results.png}
    \caption{\textbf{Latency Attack Convergence.} (Left) The Ponder Cost rises sharply from the baseline ($\approx 16$) to a saturation point ($\approx 92$) within 75 iterations. (Right) This results in a $50\%$ increase in wall-clock latency.}
    \label{fig:latency_res}
\end{figure}

\begin{figure}[h]
    \centering
    \includegraphics[width=\columnwidth]{act_pgd_energy_results.png}
    \caption{\textbf{Energy Impact.} The attack increases total energy consumption by nearly 60\% (Right) while the instantaneous power draw (Watts) remains constant, confirming a time-domain sponge effect.}
    \label{fig:energy_res}
\end{figure}

\subsection{XAI Analysis: Availability vs. Accuracy}
To understand the mechanism of the attack, we analyzed the feature attributions using Integrated Gradients (IG). Figure \ref{fig:xai_analysis} reveals a critical distinction between features important for \textit{accuracy} and those important for \textit{availability}.

The "Feature Importance" chart (Bottom Right) shows that while the \textbf{Power} feature has the highest IG Attribution (blue bar)—meaning it is critical for the prediction—the attack focuses its perturbation energy (orange bars) on environmental features like \textbf{Temperature}, \textbf{Wind Direction}, and \textbf{Wind Gusts}.
This suggests that while the model looks at previous power output to predict the next value, it relies on complex environmental variables to determine "how hard to think." The attacker exploits this by creating "chaotic weather" scenarios (high gusts/variable temperature) that maximize internal uncertainty.

\begin{figure}[h]
    \centering
    \includegraphics[width=\columnwidth]{act_pgd_latency_xai_analysis.png}
    \caption{\textbf{XAI-Guided Attack Analysis.} (Bottom Right) A key insight: While 'Power' (blue bar) is the most important feature for prediction, the attacker targets 'Wind Gusts' and 'Temperature' (orange bars) to cause latency. The model 'thinks' hardest during chaotic weather.}
    \label{fig:xai_analysis}
\end{figure}

\subsection{Attack Transferability and Defense}
Finally, we evaluated whether sponge examples generated for the ACT-LSTM affect static architectures. Figure \ref{fig:transferability} shows the cross-model performance.
\begin{itemize}
    \item \textbf{ACT-LSTM:} Sufferers massive spikes (Blue bars) under its own adversarial examples.
    \item \textbf{DeepAR (Static):} Remains completely unaffected (Orange bars), with latency staying flat at $\approx 10$ ms even when fed the "Energy Adversarial" inputs.
\end{itemize}
This confirms that \textbf{Sponge Attacks are architecture-specific}. The "loops" exploited in the ACT model do not exist in the static DeepAR graph. This empirically validates that **Architectural Distillation** (replacing dynamic models with static ones) is a perfect defense against availability attacks.

\begin{figure}[h]
    \centering
    \includegraphics[width=\columnwidth]{attack_transferability_analysis.png}
    \caption{\textbf{Transferability Analysis.} Attacks optimized for the dynamic ACT-LSTM (Blue) cause significant latency spikes on the target but have zero effect on the static DeepAR model (Orange), confirming that static architectures are naturally robust.}
    \label{fig:transferability}
\end{figure}

\section{Discussion: Mechanism of Action}
\label{sec:discussion}

While the experimental results quantify the impact of the sponge attack, understanding its underlying mechanism is crucial for designing robust defenses. In this section, we analyze the spectral properties of the adversarial perturbations and the correlation between input features and computational depth.

\subsection{Spectral Stealthiness}
A primary constraint for adversarial attacks in the Smart Grid is stealth. Large, perceptible changes to sensor readings (e.g., wind speed jumping from 5 m/s to 50 m/s) would be immediately flagged by outlier detectors.

We performed a Fast Fourier Transform (FFT) analysis of the generated perturbations $\delta = x_{adv} - x_{benign}$. As shown in Figure \ref{fig:frequency_analysis} (Top Right), the PGD attack concentrates its energy in the high-frequency domain.
\begin{itemize}
    \item \textbf{High-Frequency Injection:} The perturbation spectrum shows distinct peaks at frequencies $f > 0.3$ Hz. In the time domain (Top Left), this manifests as rapid, low-amplitude oscillations rather than global shifts.
    \item \textbf{Sensor Noise Mimicry:} This high-frequency pattern statistically resembles thermal sensor noise or quantization error, making it difficult to distinguish from valid instrument variance using standard low-pass filters.
\end{itemize}

\begin{figure}[t]
    \centering
    \includegraphics[width=\columnwidth]{frequency_analysis.png}
    \caption{\textbf{Spectral Analysis of Perturbations.} (Top) The attack introduces high-frequency oscillations that mimic sensor noise. (Bottom) The spectral power distribution shows that the attack injects the most energy into the 'Temperature' and 'Dewpoint' channels to trigger latency.}
    \label{fig:frequency_analysis}
\end{figure}

\subsection{Feature-Latency Correlation}
To determine which specific input features drive the ACT-LSTM's halting unit, we computed the Pearson correlation coefficient ($r$) between the feature values and the resulting change in ponder steps ($\Delta N$).

Figure \ref{fig:correlations} presents the results:
\begin{itemize}
    \item \textbf{Power ($r \approx +0.40$):} There is a strong positive correlation between the historical Power output and latency. The model "thinks harder" when the grid load is high, likely because the penalty for prediction error is higher in high-load regimes.
    \item \textbf{Wind Speed ($r \approx -0.35$):} Interestingly, Wind Speed shows a negative correlation. This implies that the model halts quickly when wind speeds are high (predictable generation) but hesitates when wind speeds are low or variable (intermittency).
\end{itemize}
The attacker exploits this learned bias by injecting perturbations that artificially lower the perceived wind speed while increasing the perceived load, forcing the model into its "uncertainty zone."

\begin{figure}[h]
    \centering
    \includegraphics[width=\columnwidth]{act_ponder_correlation.png}
    \caption{\textbf{Mechanism of Latency.} (Right) Correlation analysis reveals that the Halting Unit is most sensitive to the 'Power' feature ($r=0.4$). (Left) Scatter plot showing that while the attack is stochastic, perturbations generally push the model towards a specific 'uncertainty manifold' where ponder depth is maximized.}
    \label{fig:correlations}
\end{figure}

\subsection{The "Logic Trap" Hypothesis}
Combining these findings, we formulate the \textit{Logic Trap Hypothesis}: Dynamic neural networks learn a latent "complexity manifold"—a region of the input space where the prediction is difficult (e.g., sudden weather shifts). The Sponge Attack functions by computing the gradient towards this manifold. Unlike standard integrity attacks that push inputs across a linear decision boundary, sponge attacks push inputs towards the center of high-curvature regions in the loss landscape, trapping the recurrent loop in an extended state of refinement.

\subsection{Economic Impact Analysis}
Beyond operational stability, sponge attacks impose a tangible financial cost on the operator. We estimate this cost using the "Energy-per-Token" metric.

\subsubsection{Cloud Computing Costs}
For a utility processing $10^6$ inference requests per hour on AWS \textit{p3.2xlarge} instances:
\begin{itemize}
    \item \textbf{Benign Cost:} At 47ms/request, the compute cost is approximately \$11.50 per hour.
    \item \textbf{Attack Cost:} Under the sponge attack (70ms/request), the required compute capacity increases by $48\%$. This forces the auto-scaler to provision additional instances, raising the cost to $\approx \$17.02$ per hour.
    \item \textbf{Annualized Loss:} Over a year, this stealthy latency inflation results in an excess expenditure of $\approx \$48,000$ purely in wasted cloud cycles, not accounting for potential regulatory fines due to delayed reporting.
\end{itemize}

\subsubsection{Battery Replacement Costs}
For edge devices (e.g., IoT sensors with 3000mAh batteries), the 60\% increase in energy consumption (Fig. \ref{fig:energy_res}) reduces the device lifespan from 5 years to 3.1 years. This nearly doubles the replacement frequency (maintenance truck rolls), which is the dominant cost factor in AMI operations.

\section{Defense Evaluation and Trade-offs}
\label{sec:mitigation}

Defending against availability attacks in the Smart Grid requires a fundamental shift from optimizing for \textit{accuracy} to optimizing for \textit{predictability}. Unlike integrity attacks, where the goal is to detect deviations in the output space $Y$, sponge defenses must monitor the computational path in the execution space $Z$. We formally evaluate three defense paradigms, quantifying the \textit{Accuracy-Availability Trade-off}.

\subsection{Defense A: Deterministic Compute-Aware Timeouts}
The most direct defense against dynamic execution attacks is to enforce a hard upper bound on inference time, effectively converting a dynamic model into a static one under stress.

\subsubsection{Formalization}
Let $\tau_{max}$ be the maximum allowable dispatch latency (e.g., 50ms). We define a "Safe Inference" function $f_{safe}(x)$ that interrupts the recurrence loop:
\begin{equation}
    f_{safe}(x) = \begin{cases} 
        f_\theta(x) & \text{if } t_{exec} < \tau_{max} \\ 
        \text{Proj}_{\mathcal{Y}}(h_{t_{now}}) & \text{if } t_{exec} \ge \tau_{max} 
    \end{cases}
\end{equation}
where $\text{Proj}_{\mathcal{Y}}(h_{t_{now}})$ projects the current intermediate hidden state $h$ to a valid output prediction.

\subsubsection{Evaluation}
While this guarantees availability ($P(\text{DoS}) = 0$), it introduces an \textit{Accuracy Cost}. On our ACT-LSTM, enforcing $\tau_{max}=50ms$ clipped the computation of 1.2\% of benign "hard" samples (e.g., rapid storm fronts), increasing the Mean Squared Error (MSE) from 0.042 to 0.048. However, against the PGD Sponge Attack, it neutralized the latency spike entirely, reducing the worst-case time from 70.69ms to 50.00ms.

\subsection{Defense B: Architectural Distillation}
To remove the attack surface entirely, we propose \textit{Student-Teacher Distillation}. The goal is to compress the knowledge of the dynamic ACT-LSTM (Teacher $T$) into a static CNN-LSTM (Student $S$) that lacks control-flow loops.

\subsubsection{Distillation Objective}
We train the Student $S$ to mimic the probability distribution of the Teacher using a composite loss function:
\begin{equation}
    \mathcal{L}_{distill} = \alpha \mathcal{L}_{MSE}(S(x), y_{true}) + (1-\alpha) \mathcal{L}_{KL}(S(x) || T(x))
\end{equation}
where $\mathcal{L}_{KL}$ is the Kullback-Leibler divergence. 

\subsubsection{Results}
As shown in Table \ref{tab:defense_comparison}, the distilled static model achieved 96\% of the Teacher's accuracy but exhibited \textbf{zero latency variance} under attack. This confirms that static graphs are the only mathematically provable defense against algorithmic complexity attacks.

\subsection{Defense C: Energy-Regularized Adversarial Training}
For scenarios where dynamic architectures are strictly required, we propose \textit{Energy-Aware Adversarial Training (EA-AT)}. Unlike standard AT which minimizes error on perturbed inputs ($\max_\delta \mathcal{L}_{err}$), EA-AT adds a regularization term for computational cost.

\begin{equation}
    \min_\theta \mathbb{E}_{(x,y) \sim \mathcal{D}} \left[ \max_{\delta \in \Delta} \left( \mathcal{L}_{task}(f_\theta(x+\delta), y) + \lambda \Omega(N(x+\delta)) \right) \right]
\end{equation}
where $\Omega(N)$ penalizes the number of ponder steps $N$.
\begin{itemize}
    \item \textbf{Outcome:} Training with $\lambda=0.1$ reduced the average ponder depth on adversarial examples from 7.64 to 5.20.
    \item \textbf{Limitation:} The model learned to be "lazy," reducing its ability to adapt to legitimate complex weather patterns, resulting in a 3.5\% drop in benign accuracy.
\end{itemize}

\begin{table}[h]
    \centering
    \caption{Performance Trade-offs of Mitigation Strategies.}
    \label{tab:defense_comparison}
    \renewcommand{\arraystretch}{1.2}
    \begin{tabular}{l c c c}
        \toprule
        \textbf{Defense Strategy} & \textbf{DoS Risk} & \textbf{Acc. Drop} & \textbf{Deploy Cost} \\
        \midrule
        No Defense & \textbf{Critical} & 0.0\% & Low \\
        Hard Timeouts & None & 1.2\% & Low \\
        Energy-Aware Training & Low & 3.5\% & High (Retrain) \\
        \textbf{Distillation (Recommended)} & \textbf{None} & \textbf{4.0\%} & \textbf{Medium} \\
        \bottomrule
    \end{tabular}
\end{table}

\section{Future Research Directions}
\label{sec:future}

This work established the baseline vulnerability of dynamic grid AI to inference-time sponge attacks. However, the intersection of AI safety and grid hardware presents several open challenges for future investigation.

\subsection{Federated Sponge Poisoning}
Smart Grids are increasingly adopting Federated Learning (FL) to preserve prosumer privacy. In this paradigm, edge devices (smart meters) train local models and send weight updates $\Delta w$ to a central aggregator.
A malicious prosumer could execute a \textit{Sponge Poisoning Attack} by submitting crafted updates that maximize the energy consumption of the global model. Unlike data poisoning, this attack does not degrade accuracy, making it invisible to standard malicious robust aggregators (e.g., Krum or Median). Future work should investigate "Energy-Audit Aggregation" mechanisms that verify the computational cost of weight updates before integration.

\subsection{Hardware-Specific Thermal Attacks}
Our experiments focused on algorithmic latency ($O(N)$ complexity). However, a more insidious vector is the \textit{Thermal Sponge}.
\begin{itemize}
    \item \textbf{Mechanism:} An attacker could identify inputs that trigger high-intensity matrix operations (e.g., dense AVX instructions) at the resonant frequency of the edge device's heat sink.
    \item \textbf{Impact:} This would force the device's Dynamic Voltage and Frequency Scaling (DVFS) governor to throttle the CPU frequency to prevent overheating. In a substation environment, such "thermal toggling" could not only cause latency spikes but also physically degrade the silicon lifespan of critical controllers (e.g., NVIDIA Jetsons or Raspberry Pis used in microgrids).
\end{itemize}

\section{Conclusion}
This report investigated the availability threat posed by Sponge Attacks to AI-based Smart Grid systems. We identified that while traditional static RNNs (DeepAR) and Foundation Models (Chronos-v2) exhibit inherent robustness due to their fixed computational graphs and input quantization, dynamic architectures like ACT-LSTM are critically vulnerable. Our Proof-of-Concept demonstrated that a black-box Genetic Algorithm could induce a 70\% latency spike in dynamic models by exploiting the logical decision boundaries of the halting unit.

As Smart Grid operators increasingly look towards dynamic and adaptive AI for efficiency, the attack surface for Denial-of-Service expands. We conclude that \textit{availability defenses}—such as strict compute timeouts and architectural distillation—are mandatory requirements for deploying dynamic deep learning in time-critical energy infrastructure.

\begin{thebibliography}{1}

\bibitem{shumailov2021sponge}
I. Shumailov, Y. Zhao, D. Bates, N. Papernot, R. Mullins, and R. Anderson, ``Sponge examples: Energy-latency attacks on neural networks,'' in \emph{2021 IEEE European Symposium on Security and Privacy (EuroS\&P)}, IEEE, 2021, pp. 212--231.

\bibitem{sanchez2024attacking}
G. Sánchez, G. Elbez, and V. Hagenmeyer, ``Attacking Learning-based Models in Smart Grids: Current Challenges and New Frontiers,'' in \emph{Proceedings of the 15th ACM International Conference on Future and Sustainable Energy Systems}, 2024.

\bibitem{cina2022sponge}
A. E. Cinà, et al., ``Energy-latency attacks via sponge poisoning,'' \emph{arXiv preprint arXiv:2203.08147}, 2022.

\bibitem{shapira2023phantom}
A. Shapira, A. Zolfi, L. Demetrio, B. Biggio, and A. Shabtai, ``Phantom Sponges: Exploiting Non-Maximum Suppression to Attack Deep Object Detectors,'' in \emph{Proceedings of the IEEE/CVF Winter Conference on Applications of Computer Vision (WACV)}, 2023.

\bibitem{telintelo2024skip}
J. te Lintelo, S. Koffas, and S. Picek, ``The SkipSponge Attack: Sponge Weight Poisoning of Deep Neural Networks,'' \emph{arXiv preprint arXiv:2402.06357}, 2024.

\bibitem{wang2023energy}
Z. Wang, Y. Huang, et al., ``Energy-Latency Attacks to On-Device Neural Networks via Sponge Poisoning,'' \emph{arXiv preprint arXiv:2305.03888}, 2023.

\bibitem{attention}
A. Vaswani et al., ``Attention is all you need,'' in \emph{Advances in Neural Information Processing Systems}, 2017, pp. 5998--6008.

\bibitem{steganographic2026}
G. Sánchez, M. Qasim, G. Elbez, and V. Hagenmeyer, ``Steganographic Data Exfiltration for Model Stealing: A Case Study on Energy Critical Infrastructure IEC 61850 Datasets,'' in \emph{IEEE Big Data}, 2026.

\bibitem{hong2024sponge}
J. Hong, et al., ``Sponge Attack Against Multi-Exit Networks with Data Poisoning,'' \emph{IEEE Access}, vol. 12, 2024.

\bibitem{telintelo2025skipsponge}
J. te Lintelo, S. Koffas, and S. Picek, ``The SkipSponge Attack: Sponge Weight Poisoning of Deep Neural Networks,'' \emph{ITU Journal on Future and Evolving Technologies}, vol. 6, no. 3, 2025.

\bibitem{hasan2025sensing}
S. M. Hasan, H. Zangoti, et al., ``Sponge Attacks on Sensing AI: Energy-Latency Vulnerabilities and Defense via Model Pruning,'' \emph{arXiv preprint arXiv:2505.06454}, 2025.

\bibitem{sanchez2025lightweight}
G. Sánchez, G. Elbez, and V. Hagenmeyer, ``Lightweight Moving Target Defense for Robust Intrusion Detection in Smart Grids,'' in \emph{Proceedings of the IEEE International Conference on Communications, Control, and Computing Technologies for Smart Grids (SmartGridComm)}, 2025.

\end{thebibliography}

\end{document}